\begin{proposition}[determinant of the printed period derivative]
\label{prop:period-derivative-determinant}
\Verified\ With dots denoting differentiation in the coordinate on
$\mathfrak h_z$, set
\[
 K=\det(\dot Z)=6\dot\mu^{\,2}-\dot\tau\dot\beta.
\]
Then $K$ is not identically zero, so the Kodaira--Spencer map has rank two
on a dense open set.  However $K(z_1)=0$ at the chosen order-three elliptic
point.  In particular, the assertion $\dot Z\ne0$ does not imply that
$\dot Z$ is invertible, and direct images across the finite fibres require
their actual cyclic invariant lattices.
\end{proposition}

\begin{proof}
At the cusp, \cref{prop:period-input-audit} gives
\[
 t=e^{2\pi i s},\qquad Z=sB_0+C(t),\qquad
 B_0=\begin{pmatrix}0&1\\-1&0\end{pmatrix},
\]
with $C$ holomorphic at $t=0$.  Hence
\[
 \frac{dZ}{ds}=B_0+2\pi i\,tC'(t),
 \qquad
 \det\!\left(\frac{dZ}{ds}\right)=1+O(t).
\]
The cusp coordinate $s$ is locally biholomorphic on the marked cover, so
$K$ is nonzero near the punctured cusp and therefore is not identically
zero.

Now put $\rho=e^{\pi i/3}=\tau(z_1)$.  The difference between $\mu$ and
the local affine solution
\[
 \mu_1=\frac{2-\tau}{3}
\]
vanishes to order at least two at $z_1$: the homogeneous line has local
generator of order two there.  Since $\tau-\rho$ has a simple zero,
\[
 \dot\mu(z_1)=-\frac13\dot\tau(z_1).
\]
Write
\[
 \phi_1=2-\frac{6(1-\mu)^2}{\tau},
 \qquad \beta\circ g_1=\beta+\phi_1.
\]
Using $\rho^2-\rho+1=0$ and the preceding derivative gives
\[
 \dot\phi_1(z_1)
 =-\frac23(1+\rho)\dot\tau(z_1).
\]
The derivative of the linearized order-three action at its fixed point is
$g_1'(z_1)=-\rho$.  Differentiating the transformation law for $\beta$
therefore yields
\[
 -(1+\rho)\dot\beta(z_1)=\dot\phi_1(z_1),
 \qquad
 \dot\beta(z_1)=\frac23\dot\tau(z_1).
\]
Substitution gives
\[
 K(z_1)=6\left(-\frac13\dot\tau(z_1)\right)^2
 -\dot\tau(z_1)\left(\frac23\dot\tau(z_1)\right)=0.
\]
\end{proof}
